\documentclass[12pt, oneside]{book}

% Tamaño Smartphone (9:16)
\usepackage[
    paperwidth=108mm,
    paperheight=192mm,
    margin=8mm,
    bottom=15mm,
    footskip=8mm
]{geometry}

% preambulo_base.tex
\usepackage[spanish, es-tabla]{babel}
\usepackage{mathtools}
\usepackage{amsthm}
\usepackage{graphicx}
\usepackage{xcolor}

% --- FUENTES ---
\usepackage{unicode-math}
\usepackage{libertinus-otf} % Carga Libertinus Serif, Sans y Math automáticamente

% Fuente para código Lean (JuliaMono)
\setmonofont{JuliaMono-Regular.ttf}[
    Path = /home/jalonso/.local/share/fonts/juliamono/,
    Scale = MatchLowercase,
    BoldFont = JuliaMono-Bold.ttf,
    ItalicFont = JuliaMono-RegularItalic.ttf,
    BoldItalicFont = JuliaMono-BoldItalic.ttf
]

% --- CONFIGURACIÓN DE MINTED (Lean) ---
\usepackage{minted}
\renewcommand{\MintedPygmentize}{/home/jalonso/.local/bin/pygmentize}
\usemintedstyle{emacs}
\setminted{
    breaklines=true,
    frame=lines,
    framesep=2mm,
    baselinestretch=1,
    linenos=false
}

\setcounter{secnumdepth}{3}
\setcounter{tocdepth}{3}

%%% Local Variables:
%%% mode: LaTeX
%%% TeX-master: t
%%% End:


% Ajustes específicos para Móvil
\setminted{fontsize=\footnotesize} % Más pequeño para evitar cortes
\setlength{\parskip}{1.2ex}
\setlength{\parindent}{0pt}         % Sin sangría para ganar ancho
\linespread{1.1}                    % Más aire para lectura táctil

% Microtype es vital en columnas estrechas para evitar huecos blancos
\usepackage{microtype}

% Títulos compactos para no desperdiciar pantalla
\usepackage{titlesec}
\titleformat{\chapter}[hang]{\huge\bfseries\color{blue!40!black}}{\thechapter. }{0pt}{\huge}
\titlespacing*{\chapter}{0pt}{-15pt}{10pt}

% Cabeceras simplificadas
\usepackage{fancyhdr}
\pagestyle{fancy}
\fancyhead{}
\fancyhead[L]{\small\sffamily\color{gray} \nouppercase{\rightmark}}
\fancyhead[R]{\small\sffamily \thepage}
\cfoot{}

\usepackage[
  colorlinks=true,
  linkcolor=blue!70!black,
  urlcolor=blue,
  pdfauthor={Jose A. Alonso},
  pdftitle={Análisis en Lean 4 (Móvil)},
  pdfstartview=FitV]{hyperref}

\input definiciones
\input contenido

%%% Local Variables:
%%% mode: LaTeX
%%% TeX-master: t
%%% End:
